\documentclass[tikz,border=10pt]{standalone}
\usepackage[T5]{fontenc}
\usepackage[utf8]{inputenc}
\usepackage{amsmath}
\usepackage{amsfonts}
\usepackage{amssymb}
\usetikzlibrary{arrows.meta, calc, positioning}

\begin{document}
\begin{tikzpicture}[scale=1.3, >=Stealth, font=\small]

    % ==========================================
    % BÊN TRÁI: CHIẾN LƯỢC ONE-VS-ALL (OVA / OvR)
    % ==========================================
    \node[align=center, font=\bfseries\large] at (-2.5, 3.2) {Chiến lược One-vs-All (OVA) \\ $k=3$ bộ phân loại};

    % Khung tọa độ
    \draw[thick, gray!30] (-4.5, -2.5) rectangle (-0.5, 2.5);

    % Data Points (3 Lớp A, B, C)
    % Lớp A (Blue)
    \foreach \p in {(-2.5, 1.8), (-2.2, 1.5), (-2.8, 1.6), (-2.4, 1.2), (-2.6, 2.1)} {
        \fill[blue] \p circle (2pt);
    }
    \node[font=\bfseries, text=blue] at (-2.5, 2.2) {Lớp A};

    % Lớp B (Red)
    \foreach \p in {(-3.8, -1.2), (-3.5, -1.5), (-4.1, -1.0), (-3.6, -0.8), (-3.9, -1.7)} {
        \fill[red] \p circle (2pt);
    }
    \node[font=\bfseries, text=red] at (-3.8, -2.0) {Lớp B};

    % Lớp C (Green)
    \foreach \p in {(-1.2, -1.2), (-1.5, -1.5), (-0.9, -1.0), (-1.4, -0.8), (-1.1, -1.7)} {
        \fill[green!70!black] \p circle (2pt);
    }
    \node[font=\bfseries, text=green!70!black] at (-1.2, -2.0) {Lớp C};

    % Ranh giới OVA (Mỗi lớp vs Phần còn lại)
    % Ranh giới A vs (B+C)
    \draw[very thick, blue, dashed] (-4.2, 0.5) -- (-0.8, 0.5) node[right, font=\footnotesize] {A vs (B,C)};
    
    % Ranh giới B vs (A+C)
    \draw[very thick, red, dashed] (-2.5, -2.2) -- (-2.8, 1.0) node[above, font=\footnotesize] {B vs (A,C)};
    
    % Ranh giới C vs (A+B)
    \draw[very thick, green!70!black, dashed] (-2.5, -2.2) -- (-2.2, 1.0) node[above, font=\footnotesize] {C vs (A,B)};


    % ==========================================
    % BÊN PHẢI: CHIẾN LƯỢC ONE-VS-ONE (OVO)
    % ==========================================
    \node[align=center, font=\bfseries\large] at (2.5, 3.2) {Chiến lược One-vs-One (OVO) \\ $k(k-1)/2 = 3$ bộ phân loại};

    % Khung tọa độ
    \draw[thick, gray!30] (0.5, -2.5) rectangle (4.5, 2.5);

    % Data Points (3 Lớp A, B, C) - Vị trí tương tự bên trái
    % Lớp A (Blue)
    \foreach \p in {(2.5, 1.8), (2.2, 1.5), (2.8, 1.6), (2.4, 1.2), (2.6, 2.1)} {
        \fill[blue] \p circle (2pt);
    }
    \node[font=\bfseries, text=blue] at (2.5, 2.2) {Lớp A};

    % Lớp B (Red)
    \foreach \p in {(1.2, -1.2), (1.5, -1.5), (0.9, -1.0), (1.4, -0.8), (1.1, -1.7)} {
        \fill[red] \p circle (2pt);
    }
    \node[font=\bfseries, text=red] at (1.2, -2.0) {Lớp B};

    % Lớp C (Green)
    \foreach \p in {(3.8, -1.2), (3.5, -1.5), (4.1, -1.0), (3.6, -0.8), (3.9, -1.7)} {
        \fill[green!70!black] \p circle (2pt);
    }
    \node[font=\bfseries, text=green!70!black] at (3.8, -2.0) {Lớp C};

    % Ranh giới OVO (Từng cặp Lớp i vs Lớp j)
    % Ranh giới A vs B
    \draw[very thick, purple, dashed] (0.8, 0.8) -- (3.2, -0.5) node[right, font=\footnotesize] {A vs B};
    
    % Ranh giới A vs C
    \draw[very thick, orange!90!black, dashed] (4.2, 0.8) -- (1.8, -0.5) node[left, font=\footnotesize] {A vs C};
    
    % Ranh giới B vs C
    \draw[very thick, teal, dashed] (2.5, -2.3) -- (2.5, 0.5) node[above, font=\footnotesize] {B vs C};


    % --- Caption ---
    \node[draw, fill=white, thick, align=center] at (0, -3.2) {
        \textbf{Hình 1.4.1 \& 1.4.2: Đối sánh ranh giới phân chia giữa OVA và OVO} \\
        OVA tìm ranh giới phân biệt 1 lớp với tất cả các lớp còn lại (dễ mất cân bằng dữ liệu). \\
        OVO tìm ranh giới phân biệt riêng từng cặp lớp (dữ liệu cân bằng hơn nhưng số lượng mô hình tăng theo $O(k^2)$).
    };

\end{tikzpicture}
\end{document}
